\section{Linear Code}
%We cover the following topics in this section.
%\begin{itemize}
%	\item Finite fields
%	\item Linear encodings
%	\item \([n,k]\) linear codes
%	\item Generator matrices, systematic form, repetition
%	\item Parity-check codes, and parity-check matrices.
%\end{itemize}
%\hrule\vspace{12pt}
\definecolor{accent}{HTML}{1F4E79}
\definecolor{softaccent}{HTML}{EAF2F8}
\definecolor{softgray}{HTML}{F6F7F9}
\definecolor{successdraw}{HTML}{27AE60}
\definecolor{successfill}{HTML}{EAFAF1}
\definecolor{errordraw}{HTML}{C0392B}
\definecolor{errorfill}{HTML}{FDEDEC}
%\begin{center}
%\begin{tikzpicture}[scale=1,>=Stealth,
%	block/.style={
%		draw=accent,
%		thick,
%		fill=softaccent,
%		rounded corners=2mm,
%		minimum width=3.2cm,
%		minimum height=1cm,
%		align=center,
%		font=\sffamily,
%		drop shadow={opacity=0.15, shadow xshift=1mm, shadow yshift=-1mm} % Subtle depth
%	},
%	arrow/.style={
%		->,
%		thick,
%		accent,
%		rounded corners=2mm % Softens the sharp 90-degree turns
%	}]
%	\node[block] (tx) at (0,0) {Data $\vect{x}$};
%	\node[block] (enc) at (6,0) {Encoder};
%	\node[block] (chan) at (10,-2) {Channel + noise};
%	\node[block] (dec) at (6,-4) {Decoder};
%	\node[block] (rx) at (0,-4) {Data $\hat{\vect{x}}$};
%	\draw[arrow] (tx) -- node[above] {$\vect{x}\in\F_q^k$} (enc);
%	\draw[arrow] (enc) -| node[above] {$\vect{c}=\vect{x}G\in\F_q^n$} (chan);
%	\draw[arrow] (chan) |- node[below] {$\vect{y}=\vect{c}+\vect{e}\in\F_q^n$} (dec);
%	\draw[arrow] (dec) -- (rx);
%\end{tikzpicture}
%\end{center}
\begin{center}
\begin{tikzpicture}[scale=1,>=Stealth,
	block/.style={
		draw=accent,
		thick,
		fill=softaccent,
		rounded corners=2mm,
		minimum width=3.2cm,
		minimum height=1cm,
		align=center,
		font=\sffamily,
		drop shadow={opacity=0.15, shadow xshift=1mm, shadow yshift=-1mm}
	},
	decision/.style={
		diamond,
		aspect=1.5, % Stretches the diamond horizontally to fit text
		draw=accent,
		thick,
		fill=softgray,
		align=center,
		font=\sffamily\small,
		drop shadow={opacity=0.15, shadow xshift=1mm, shadow yshift=-1mm}
	},
	successblock/.style={
		block,
		draw=successdraw,
		fill=successfill,
		text=black!85
	},
	errorblock/.style={
		block,
		draw=errordraw,
		fill=errorfill,
		text=black!85
	},
	arrow/.style={
		->,
		thick,
		accent,
		rounded corners=2mm 
	},
	dashedarrow/.style={
		->,
		thick,
		accent,
		dashed, % Distinguishes evaluation logic from data flow
		rounded corners=2mm 
	}]
	
	% --- Main Pipeline Nodes ---
	\node[block] (tx) at (0,0) {Data $\vect{x}$};
	\node[block] (enc) at (6,0) {Encoder};
	\node[block] (chan) at (10,-1.5) {Channel + noise};
	\node[block] (dec) at (6,-3) {Decoder};
	\node[block] (rx) at (0,-3) {Data $\hat{\vect{x}}$};
	
	% --- Evaluation Nodes ---
	\node[decision] (comp) at (0,-6) {Is $\hat{\vect{x}} = \vect{x}?$};
	\node[successblock] (succ) at (6,-5) {\textbf{Successful Decoding}\\$\hat{\vect{x}} = \vect{x}$};
	\node[errorblock] (err) at (6,-7) {\textbf{Decoding Error}\\$\hat{\vect{x}} \neq \vect{x}$};
	
	% --- Main Pipeline Arrows ---
	\draw[arrow] (tx) -- node[above, font=\small] {$\vect{x}\in\F_q^k$} (enc);
	\draw[arrow] (enc) -| node[above right, font=\small, pos=0.25] {$\vect{c}=\vect{x}G\in\F_q^n$} (chan);
	\draw[arrow] (chan) |- node[below right, font=\small, pos=0.75] {$\vect{y}=\vect{c}+\vect{e}\in\F_q^n$} (dec);
	\draw[arrow] (dec) -- (rx);
	
	% --- Evaluation Arrows ---
	% Route from estimated data down to comparison
	\draw[dashedarrow] (rx.south) -- node[right, font=\small] {Estimate} (comp.north);
	
	% Route from original data around the left side down to comparison
	\draw[dashedarrow] (tx.west) -- node[above, font=\small] {Original} ++(-1.5,0) |- (comp.west);
	
	% Branching arrows from comparison to success/error
	\draw[arrow] (comp.east) -- ++(1.5,0) |- (succ.west);
	\draw[arrow] (comp.east) -- ++(1.5,0) |- (err.west);
	
\end{tikzpicture}
\end{center}
\noindent
A generic communication model has the form
\[
\vect{x}=(x_1,\dots,x_k)
\xrightarrow{\text{encoder}}
\underbrace{\vect{c}=(c_1,\dots,c_n)}_{\text{codeword}\; (n\geq k)}
\xrightarrow{\text{noise }\vect e}
\vect{y}=\vect{c}+\vect{e}
\xrightarrow{\text{decoder}}
\hat{\vect{x}}=(\hat x_1,\dots,\hat x_k).
\]

\newpage

\begin{example}[Repetition code]
Let \(\Sigma\) be an alphabet. An \emph{encoding} is a map
\[
\fullfunction{\mathsf{Encode}}{\Sigma^k}{\Sigma^n}{(x_1,\dots,x_k)}{(c_1,\dots,c_n)}
\] with \(n\ge k\). %The output vector \(\vect c=(c_1,\dots,c_n)\) is the transmitted codeword.
When \(k=1\), a very simple encoding repeats the same symbol
%\[
%(x_1)\longmapsto (x_1,x_1,\dots,x_1)\in \Sigma^n.
%\]
\begin{center}
	\begin{tikzpicture}[
		>=Stealth,
		% Style for single input symbols
		symbol/.style={
			circle, draw=accent, thick, fill=softgray,
			minimum size=9mm, inner sep=0pt, align=center, font=\normalsize,
			drop shadow={opacity=0.15, shadow xshift=1mm, shadow yshift=-1mm}
		},
		% Style for the output codeword arrays
		block/.style={
			rectangle, draw=accent, thick, fill=softaccent, rounded corners=2mm,
			minimum width=3.8cm, minimum height=1.1cm, align=center, font=\normalsize,
			drop shadow={opacity=0.15, shadow xshift=1mm, shadow yshift=-1mm}
		},
		% Style for the mapping arrows
		maparrow/.style={
			|->, thick, accent, shorten >=1.5mm, shorten <=1.5mm
		},
		labeltext/.style={
			font=\small\sffamily, text=black!85
		}
		]
		% --- General Formula Row ---
		\node[symbol] (x1) at (0, 0) {$x_1$};
		\node[block] (outX) at (5.5, 0) {$(x_1, \, x_1, \, \dots, \, x_1)$};
		
		\draw[maparrow] (x1) -- node[above, labeltext] {$\mathsf{Encode}$} (outX.west);
		
		% Brace under the general codeword to show length
		\draw[thick, accent, decorate, decoration={brace, amplitude=5pt, mirror}]
		([xshift=4mm, yshift=-2.5mm]outX.south west) -- 
		([xshift=-4mm, yshift=-2.5mm]outX.south east)
		node[midway, below=6pt, labeltext] {$n$ copies ($\in \Sigma^n$)};
	\end{tikzpicture}
\end{center}
%If \(\Sigma=\{0,1\}\) then
%\[
%(0)\mapsto (0,0,\dots,0),
%\qquad
%(1)\mapsto (1,1,\dots,1).
%\]
\begin{center}
	\begin{tikzpicture}[
		>=Stealth,
		% Style for single input symbols
		symbol/.style={
			circle, draw=accent, thick, fill=softgray,
			minimum size=9mm, inner sep=0pt, align=center, font=\normalsize,
			drop shadow={opacity=0.15, shadow xshift=1mm, shadow yshift=-1mm}
		},
		% Style for the output codeword arrays
		block/.style={
			rectangle, draw=accent, thick, fill=softaccent, rounded corners=2mm,
			minimum width=3.8cm, minimum height=1.1cm, align=center, font=\normalsize,
			drop shadow={opacity=0.15, shadow xshift=1mm, shadow yshift=-1mm}
		},
		% Style for the mapping arrows
		maparrow/.style={
			|->, thick, accent, shorten >=1.5mm, shorten <=1.5mm
		},
		labeltext/.style={
			font=\small\sffamily, text=black!85
		},
		header/.style={
			font=\bfseries\sffamily, text=accent
		}
		]
		% --- Binary Case: Input 0 ---
		% Using white fill to contrast with the abstract case
		\node[symbol, fill=white] (b0) at (0, -2.8) {$0$};
		\node[block, fill=white] (out0) at (5.5, -2.8) {$(0, \, 0, \, \dots, \, 0)$};
		
		\draw[maparrow] (b0) -- (out0.west);
		
		% --- Binary Case: Input 1 ---
		\node[symbol, fill=white] (b1) at (0, -4.3) {$1$};
		\node[block, fill=white] (out1) at (5.5, -4.3) {$(1, \, 1, \, \dots, \, 1)$};
		
		\draw[maparrow] (b1) -- (out1.west);
	\end{tikzpicture}
\end{center}
This is called the \emph{repetition code}. If one symbol is changed during transmission, we can still recover the original message by \emph{majority vote}. For instance,
\[
(0,0,0)\rightsquigarrow (0,1,0),
\qquad
(1,1,1)\rightsquigarrow (1,0,1).
\]
Since \(0\) appears most often in \((0,1,0)\), we decode it as \(0\), and since
\(1\) appears most often in \((1,0,1)\), we decode it as \(1\).
%More generally, if \(n\) is odd, then the repetition code can correct up to
%\[
%\left\lfloor \frac{n-1}{2}\right\rfloor
%\]
%errors, because the original symbol is still the majority symbol.

\medskip

\begin{center}
\definecolor{success}{HTML}{27AE60}
\begin{tikzpicture}[
	>=Stealth,
	block/.style={
		rectangle, draw=accent, thick, fill=softaccent, rounded corners=2mm,
		minimum width=2.2cm, minimum height=1cm, align=center, font=\large,
		drop shadow={opacity=0.15, shadow xshift=1mm, shadow yshift=-1mm}
	},
	decision/.style={
		rectangle, draw=accent, thick, fill=softgray, rounded corners=2mm,
		minimum width=3.2cm, minimum height=1.1cm, align=center, font=\sffamily\small,
		drop shadow={opacity=0.15, shadow xshift=1mm, shadow yshift=-1mm}
	},
	result/.style={
		circle, draw=success, thick, fill=success!10, 
		minimum size=10mm, inner sep=0pt, align=center, font=\Large\bfseries, text=success,
		drop shadow={opacity=0.15, shadow xshift=1mm, shadow yshift=-1mm}
	},
	arrow/.style={
		->, thick, accent
	},
	noisearrow/.style={
		->, thick, errordraw, 
		% The snake decoration mimics the \rightsquigarrow command
		decorate, decoration={snake, amplitude=0.4mm, segment length=3mm, post length=1.5mm}
	},
	labeltext/.style={font=\small\sffamily, text=black!85},
	header/.style={font=\bfseries\sffamily, text=accent}
	]
	
	% --- Column Headers ---
	\node[header] at (0, 1.5) {Original Data};
	\node[header] at (4, 1.5) {Received};
	\node[header] at (8.5, 1.5) {Decoder};
	\node[header] at (12.8, 1.5) {Recovered Data};
	
	\node[block] (tx1) at (0, 0) {$0$};
	
	% The corrupted bit is highlighted in bold red
	\node[block] (rx1) at (4, 0) {$(0, \textcolor{errordraw}{\mathbf{1}}, 0)$};
	
	\node[decision] (vote1) at (8.5, 0) {Two 0s, One 1\\[1mm] (\textbf{0 is the majority})};
	\node[result] (out1) at (12.8, 0) {0};
	
	\draw[noisearrow] (tx1) -- node[above, font=\small\sffamily, text=errordraw] {error} (rx1);
	\draw[arrow] (rx1) -- node[above, labeltext] {} (vote1);
	\draw[arrow] (vote1) -- (out1);
	
	% --- Horizontal Divider ---
%	\draw[dashed, draw=gray!40, thick] (-1.5, -1.6) -- (14, -1.6);
	
%	% ==========================================
	% EXAMPLE 2: Sending 1s
	% ==========================================
	\node[block] (tx2) at (0, -3.2) {$1$};
	
	% The corrupted bit is highlighted in bold red
	\node[block] (rx2) at (4, -3.2) {$(1, \textcolor{errordraw}{\mathbf{0}}, 1)$};
	
	\node[decision] (vote2) at (8.5, -3.2) {Two 1s, One 0\\[1mm] (\textbf{1 is the majority})};
	\node[result] (out2) at (12.8, -3.2) {1};
	
	\draw[noisearrow] (tx2) -- node[above, font=\small\sffamily, text=errordraw] {error} (rx2);
	\draw[arrow] (rx2) -- node[above, labeltext] {} (vote2);
	\draw[arrow] (vote2) -- (out2);
\end{tikzpicture}
\end{center}
\end{example}

\newpage
\subsection{Linear Codes}
\defbox[Linear Encoding]{\begin{definition}
A \textbf{linear encoding} over \(\F_q\) is a linear map
\[
\F_q^k \longrightarrow \F_q^n,
\qquad
\vect{x}\longmapsto \vect{c}.
\]
\end{definition}}
\begin{remark}
Since the map is linear,
\begin{itemize}
	\item the sum of two codewords is again a codeword;
	\item any scalar multiple of a codeword is again a codeword.
\end{itemize}
\end{remark}

\medskip

\defbox[Linear Code]{\begin{definition}
	A linear subspace \(\code\subseteq \F_q^n\) is called a \textbf{linear code}.
\end{definition}}


If \(\code\) is a linear code, then automatically:
\begin{itemize}
	\item the zero codeword \(\vect 0\) belongs to \(\code\);
	\item whenever \(\vect c\in \code\), also \(-\vect c\in \code\);
	\item every codeword is a linear combination of basis vectors of \(\code\).
\end{itemize}

\begin{proposition}
	If \(\code\) is an \((n,k)\) linear code over \(\F_q\), then
	\[
	|\code|=q^k.
	\]
\end{proposition}

\begin{proof}
	Choose a basis \(\vect b_1,\dots,\vect b_k\) of \(\code\). Every codeword can be written uniquely as
	\[
	x_1\vect b_1+\cdots+x_k\vect b_k,
	\qquad x_i\in \F_q.
	\]
	Each coefficient \(x_i\) has \(q\) choices, and there are \(k\) independent coefficients. Therefore the number of codewords is \(q^k\).
\end{proof}
